\section{Affine Unique Games on other Globally Hypercontractive Graphs}\label{sec:other_graphs}
In this section, we explain how to adapt our techniques from Sections~\ref{sec:johnson_1},~\ref{sec:johnson_2} for 
general globally hypercontractive graphs. We first give a semi-formal definition of global hypercontractivity and go over the steps of our algorithm, noting that even though we stated it for Johnson graphs earlier, it works for general graphs as long as we have an SoS certificate of global hypercontractivity. Later we give the specifics for the proofs of Theorems~\ref{thm:main_grass} and~\ref{thm:main_hdx}: UG algorithms for the Grassmann graph and random walks over HDXs.

\subsection{The Components that Go Into the Algorithm}\label{sec:abstraction}
The $4$-components that are used in 
our algorithm proving Theorem~\ref{thm:main_johnson} 
and the interaction between them can be abstracted as follows.
\begin{enumerate}
    \item \textbf{Global hypercontractivity.}\label{def:global-hypcont} 
    The first component going into our algorithm is a global hypercontractive inequality (or a consequence of it).
    This is a result asserting that 
    for our underlying graph $G$, for all $\gamma>0$ there are $r\in\mathbb{N}$, 
    $\eps_0,\ldots,\eps_r,\zeta,\eta>0$ and a collection 
    $\mathcal{C} = \mathcal{C}_0\cup\ldots\cup\mathcal{C}_r$ of sets of vertices (these are the combinations 
    of the basic sets) satisfying the following properties: 
    \begin{enumerate}
        \item 
        \textbf{$\mathcal{C}$ consists of non-expanding sets:} 
        The edge expansion of each $C\in\mathcal{C}$ in $G$ is $1-\zeta$.
        \item 
        \textbf{$\mathcal{C}$ explains all small sets that 
        do not expand well.} A basic form of this property 
        is that sets $S\subseteq V(G)$ that are not dense in 
        any set in $\mathcal{C}$, in the sense that $\delta(S\cap C)\leq \eps_1\delta(C)$
        for all $C\in\mathcal{C}$, have expansion close to $1$, that is, $\Phi(S)\geq 1-\gamma$.
        
        A stronger version of this property, that often follows from the above (as is the case of all graphs of interest in this paper) and that is used in the proof of Theorem~\ref{thm:main_johnson_close_to_1}, is that 
        one can capture a constant fraction of $S$ using $C\in\mathcal{C}$ in which $S$ is dense. 
        Equivalently, this property says that for a set of vertices $S\subseteq V(G)$, taking $\mathcal{C}' = \{C\in\mathcal{C}~|~C\in\mathcal{C}_i,~\delta(S\cap C)\geq \eps_i \delta(C)\}$ and defining $C' = \bigcup_{C\in \mathcal{C}'}C$, if $\delta(S\cap C')\leq \eta \delta(S)$ then
        $\Phi(S)\geq 1-\gamma$.
        
        An even stronger form of global hypercontractivity, that again follows from the most basic form in our cases of interest, is the edge covering theorem. 
        We used this form of global hypercontracitivity in the proof of Theorem~\ref{thm:main_johnson}. 
        Informally, this result says that not only a substantial chunk of $S$ can be covered by non-expanding sets from $\mathcal{C}$ in which $S$ is dense, but in fact one can capture almost all of the edges staying inside $S$ using $C\in\mathcal{C}$. More precisely, 
        this result asserts that 
        we may take $\mathcal{C}''\subseteq \mathcal{C}'$ 
        such that the following properties hold:
        \begin{enumerate}
            \item Edge covering: 
        sampling an edge $(u,v)$ inside $S$, we have that 
        except with probability $\xi$ it holds that $(u,v)$ 
        is an edge inside some $C''\subseteq \mathcal{C}''$.
            \item Maximally dense: in words, this means that 
            every $C''\in\mathcal{C}''$ has no $C\supsetneq C''$ 
            in $\mathcal{C}$ in which $S$ is somewhat dense. 
            Formally, we require that for all $i=1,\ldots,r$, 
            $C''\in\mathcal{C}''\cap \mathcal{C}_i$, 
            $j<i$ and $C_j\in\mathcal{C}_j$ such that $C_{j}\supseteq C''$, it holds that 
            $\delta(S\cap C_j) < \eps_j \delta(C_j)$. 
        \end{enumerate}
    \end{enumerate}
    The edge covering theorem allows us to argue that in our shift partition
    $F_s(X,X') = \{v\in V~|~ X(v) - X'(v) = s\}$ we may capture
    almost all edges that stay within some $F_s$ using the basic 
    sets from $\mathcal{C}$. In particular, since edges that both $X$ and $X'$ satisfy constitute a constant fraction of edges and they all stay within some $F_s$, it means that we may capture almost all of them using sets from $\mathcal{C}$. 
    %Thus, we are able to find $C\in\mathcal{C}$ that becomes dense in some $F_s$ with noticeable probability over the pseudodistribution and inside $C$
    \item  By elementary arguments, it follows
    that there is such a $C\in\mathcal{C}$ and $s$ such 
    that $F_s$ 
    becomes dense in $C$ with noticeable probability and inside $C$ a constant fraction of the edges are satisfied by both $X$ and $X'$. In our argument we consider this event 
    (this was approximated by the polynomial $P_a(X,X')$ in our algorithm to stay within 
    the realm of \sos). As we explained in the introduction
    this conditioning may create correlation between $X$ and 
    $X'$ that would cause our rounding technique not to work. 
    
    To circumvent this, we showed that: 
    \begin{enumerate}
        \item \textbf{Eliminating global correlations:} using conditioning (as in~\cite{RT12}) we are able to make sure that the two collection 
    of random variables $\{Z_{u,v}\}_{u,v\in C}, \{Z'_{u,v}\}_{u,v\in C}$ of interest 
    (that are concerned with local information about the assignments $X$ and $X'$ on vertices $u,v$) are roughly
    pairwise independent.
        \item \textbf{Conditioning mostly preserves independence:} As the probability of the event that
        $C$ becomes dense in $F_s$ is noticeable, and $\{Z_{u,v}\}_{u,v\in C}$,$\{Z'_{u,v}\}_{u,v\in C}$ were pairwise-independent to begin with, we argue 
        that it cannot create too many dependencies between
        $(Z_{u,v}, Z'_{u,v})$.
    \end{enumerate} 
    
    \item \textbf{Rounding.} We then gave a rounding procedure on a subcube $C$ and a pseudodistribution $\cD$, where on average (over $\cD$) our supposed assignments $X,X'$ satisfy a constant fraction of the edges (formalized via the notion of
    shift-partition potential), some part $F_s$ of the shift partition is large, and on which we have near 
    independence between certain random variables 
    $\{(Z_{u,v},Z'_{u,v})\}_{u,v}$. The output of the rounding procedure is an assignment to the vertices of $C$ that satisfies a constant fraction of the constraints inside it.
    
    \item \textbf{An iteration result.} By that, we mean a result saying that if we have a procedure as above that is able to find a $C\in\mathcal{C}$ and find an assignment to the vertices of $C$ that satisfies a constant fraction of the constraints in it, then we can iterate it to satisfy a constant fraction 
    of the constraints in the whole graph. 
    
    This is achieved by randomizing the edges incident to $C$, which ensures that any subsequent $C'$ we find can only have negligible overlap with the $C$'s we have found thus far.
    The idea is that any $C$ that we found has at least 
    $\delta$ of the edges touching it staying inside it, 
    out of which we manage to satisfy a constant fraction, 
    hence overall we manage to satisfy a constant fraction
    of the edges touching $C$. In return for that, 
    we randomize the edges touching $C$ hence effectively giving up on the objective value coming from those edges. The point is that since the ratio between the number of edges we manage to satisfy, and the number of edges we give up on is $\Omega(1)$, after a few iterations we will manage to satisfy a 
    constant fraction of the edges before we drop the objective value of our instance by too much.
\end{enumerate}


Our of these $4$ components, the second, third and 
fourth components work as is for any of the graphs that we 
care about. Indeed, the elimination of correlations and 
conditioning are completely generic, and the rounding 
procedure only hinges on the constraints being affine. 
As for the fourth component, it only relies on the properties
of the family $\mathcal{C}$ guaranteed from our structure 
theorem and otherwise is also generic.
\skipi
We next discuss the first component, which is a globally 
hypercontractive inequality, and for concreteness we 
consider the most basic version of it. This is a result
saying that a set $S$ which is not dense inside any 
$C\in \mathcal{C}$ has large expansion, and in 
all of our graphs of interest such results are known
~\cite{KMMS,KhotMS18,filmus2020hypercontractivity,keevash2021global,bafna2022hypercontractivity,gur2022hypercontractivity}. 
For such a result to give us an algorithm, we also require 
the proof of this part to be in the \sos proof system of
constant degree (which is the case in all of the above
examples, but in principle may not be the case). 

In the case of the Johnson graph the collection 
$\mathcal{C}$ is simply the collection of basic sets 
\[
\mathcal{C} = \left\{J|_a ~|~|a|\leq r\right\},
\qquad
\text{where }
J|_a = \left\{u\in\binom{[n]}{\ell}~\Big|~u\supseteq a\right\},
\]
the basic form of the global hypercontractive inequality is 
from~\cite{KMMS} (in the appendix we give an \sos version of 
that proof), and we have shown in Section~\ref{sec:johnson_2} how to conclude from it the edge covering theorem in a black-box manner.

\subsection{The Algorithm for the Grassmann Graph}
We now explain the first component above 
for the Grassmann graph, which is established in~\cite{KhotMS18}. Recall that the Grassmann graph 
${\sf Grass}(n,\ell)$ is the graph whose vertices are all
$\ell$-dimensional subspaces $L\subseteq \mathbb{F}_2^n$, 
and $(L,L')$ is an edge if ${\sf dim}(L\cap L') = \ell-1$.


Inside the Grassmann graph ${\sf Grass}(n,\ell)$ we have the two basic sets
\[
Z_x = \left\{L~|~x\in L\right\}\text{ for $x\in \mathbb{F}_2^n$},
\qquad
Z'_W = \left\{L~|~L\subseteq W\right\}\text{ for 
a hyperplane $W\subseteq \mathbb{F}_2^n$},
\]
and using them we can define the collections $\mathcal{C}_0,\ldots,\mathcal{C}_r$ and $\mathcal{C} = \mathcal{C}_0\cup\ldots\cup\mathcal{C}_r$ as above. Indeed, 
for a subspace $X\subseteq \mathbb{F}_2^n$ we define 
$Z_X = \cap_{x\in X}Z_x$ and for a subspace $W\subseteq \mathbb{F}_2^n$ we define $Z'_W=\cap_{W'\supseteq W} Z_{W'}$ 
and then
\[
\mathcal{C}_{m} = 
\bigcup_{m_1+m_2 = m}
\left\{Z_{X}\cap Z'_{W}~\Big|{\sf dim}(X) = m_1, {\sf dim}(W) = n-m_2, X\subseteq W\right\}.
\]
The result of~\cite{KhotMS18} shows that for all $\gamma>0$, 
there are $r\in\mathbb{N}$ and $\eps>0$ 
such that if $S$ is a set of 
vertices in ${\sf Grass}(n,\ell)$ such that $\delta(S\cap C)\leq \eps \delta(C)$ for all $C\in\mathcal{C}$, then $\Phi(S)\geq 1-\gamma$. In fact, the same proof shows that a slightly more general statement holds for bounded functions, 
asserting that if $F\colon V({\sf Grass}(n,\ell))\to [0,1]$ satisfies that 
$\E_{x}[F(x)^2 1_{x\in C}]\leq \eps \E_{x}[1_{x\in C}]$, then 
\[
\langle F, T F\rangle\geq (1-\gamma)\E_x[F^2(x)].
\]
where $T$ is the normalized adjacency matrix of ${\sf Grass}(n,\ell)$. Moreover their proof can be seen 
to be \sos-able since each step uses polynomial inequalities like the Parceval, Cauchy-Schwarz and H\"{o}lder's inequality on the variables $F$. We have used precisely these inequalities in Section~\ref{sec:fourier} to get an SoS certificate of global hypercontractivity for the Johnson graph (Theorem~\ref{thm:structure-johnson}). In fact, \cite{KhotMS18}'s proof is very similar to that of~\cite{KMMS} for the Johnson graph therefore one can similarly get an SoS certificate for the Grassmann.


From this result, one can follow along the lines of Theorem~\ref{thm:edge-covering-restated} to deduce, 
in a black-box manner, an edge covering theorem for 
the Grassmann graph: 

\begin{theorem}
For all $\eta>0$ there 
are $r\in\mathbb{N}$ and
$\eps_0,\ldots,\eps_r>0$ such that 
given a set 
$S\subseteq {\sf Grass}(n,\ell)$ 
one may find $\mathcal{C}_i'\subseteq \mathcal{C}_i$ 
such that 
\begin{enumerate}
    \item For each $i$ and $C\in\mathcal{C}_i'$ it holds 
    that $\delta(S\cap C)\geq \eps_i\delta(C)$.
    \item For $\tilde{S} = \bigcup_{i=0}^{r} \bigcup_{C\in \mathcal{C}_i} C$, one has 
    \[
    \Pr_{(L,L')\in E({\sf Grass}(n,\ell))}\left[L,L'\in S\land 
    (L\not\in\tilde{S}\text{ or }L'\not\in\tilde{S})\right]\leq \eta\delta(S).
    \]
    \item For all $j<i$, $C_i\in\mathcal{C}_i'$ and 
    $C_j\in \mathcal{C}_j$ such that $C_j\supseteq C_i$ 
    it holds that $\delta(S\cap C_j)\leq \eps_j \delta(C_j)$.
\end{enumerate}
\end{theorem}
\begin{proof}[Proof sketch]
We let $\eps_0,\ldots,\eps_r$ to be determined; it will only
matter to us that $\eps_r$ is small enough. 
Take $\mathcal{C}_i'$ to be all $C_i\in \mathcal{C}_i$ 
such that $\delta(S\cap C_i)\geq \eps_i \delta(C_i)$ and 
for which there are no $j$ and $C_j\in\mathcal{C}_j$ 
such that $C_j\supseteq C_i$ and $\delta(S\cap C_j)\geq \eps_j \delta(S)$.

Let $F = 1_S$, $\tilde{F} = 1_{\tilde{S}}$, 
and note that $G = F(1-\tilde{F}) = 1_{S\cap \overline{\tilde{S}}}$ 
is $(r,\eps_r)$ pseudo-random as per the definition of~\cite{KhotMS18}, so $\Phi(S\cap \overline{\tilde{S}})\geq 1-\eta^2/100$ provided that $\eps_r$ is sufficiently small. Note that
\[
\Pr_{(L,L')\in E({\sf Grass}(n,\ell))}\left[L,L'\in S\land 
    (L\not\in\tilde{S}\text{ or }L'\not\in\tilde{S})\right]
    \leq 2\E_{(L,L')}[F[L]F[L'](1-\tilde{F}[L])],
\]
which is equal to $2\langle G, T F\rangle$, where 
$T$ is the adjacency operator of the Grassmann graph with 
self loops. 
We may bound:
\begin{align*}
\langle G, T F\rangle
=
\langle T G, F\rangle
\leq \|T G\|_2\|F\|_2
= \sqrt{\langle T G, TG\rangle} \|F\|_2
&\leq \sqrt{\langle G, TG\rangle} \|F\|_2\\
&=\|G\|_2\sqrt{1-\Phi(S\cap\overline{\tilde{S}})} \|F\|_2.
\end{align*}
We used the fact that $T$ is positive semi-definite. 
Clearly $\|G\|_2,\|F\|_2\leq \sqrt{\delta(S)}$ hence we get
$\langle G, T F\rangle\leq \delta(S)\frac{\eta}{10}$, 
and plugging that above gives the second item.
\end{proof}
Adjusting this result to the setting of the shift 
partition and phrasing it as an \sos statement 
in the standard way, one gets an analog of 
Theorem~\ref{thm:edge-covering-restated} for the 
Grassmann graph. For that, one has to address 
edges going across different $C_i$'s and 
show that they have a small contribution. 
For that one has to choose 
$\eps_{i}$ to be sufficiently smaller than 
$\eps_{i+1}$, and include in $\mathcal{C}_i'$ 
only $C_i\in\mathcal{C}_i$ that are maximally dense, and repeat a calculation analogous to the one in Theorem~\ref{thm:edge-covering-restated} when bounding the last term therein.

The rest of our algorithm then proceeds in exactly the same way.

\subsection{The algorithm for High Dimensional Expanders}\label{sec:ug-on-hdx}
In this section we consider UG instances defined over higher-order walks on two-sided local spectral expanders and discuss why our results for Johnson graphs (that can be seen as partial-swap walks over the complete complex) generalize to HDXs. We refer the reader to~\cite{dikstein2018boolean} for an excellent exposition on HDXs. 

%Any weighted undirected graph $G$ is uniquely associated with a self-adjoint random-walk matrix $M$ where $M(u,v) = \Pr_E[(u,v)]/\sum_{w}\Pr_E[(u,w)]$. We can therefore view the constraint graph in this manner instead, referring to unique games instances over random-walks as $I = (M,\Pi)$.

For concreteness, let us consider a consider a two-sided local spectral expanding complex $X$ and the partial-swap walk over $X(\ell)$ with depth $\alpha$. This random-walk naturally corresponds to a graph $G$ over the vertices $X(\ell)$. Given $G$, analogous to the subcubes in the Johnson graph, there are a family of basic sets called the links of the complex. \cite{bafna2022hypercontractivity},~\cite{gur2022hypercontractivity} generalized the result of~\cite{KMMS} to prove global hypercontractivity on HDXs and in particular showed that any non-expanding set in $G$ must have large constant density inside a link. The parameters they get are the same as that in the statement of~\cite{KMMS} for the $\alpha$-noisy Johnson graphs. In fact the proof of~\cite{bafna2022hypercontractivity} proceeded exactly along the lines of~\cite{KMMS} and therefore can be easily seen to be SoS-able to get an analogue of Theorem~\ref{thm:structure-johnson} for HDXs. One can similarly also get the stronger edge-covering theorem for HDXs (analogous to Theorem~\ref{thm:edge-covering-restated}) and we omit the details here. 


Fortunately for us,~\cite{DBLP:conf/soda/BafnaHKL22} generalized the algorithm of~\cite{BBKSS} to get a UG algorithm for HDXs albeit suffered the same two drawbacks as~\cite{BBKSS} -- an $\ell$-dependent soundness guarantee and reliance on completeness being close to $1$.~\cite{DBLP:conf/soda/BafnaHKL22} thus show the generality of the~\cite{BBKSS} framework and illustrate that it does not depend on the symmetry properties/regularity of Johnson graphs, since higher-order walks could be over irregular sparse and highly asymmetric graphs. We can similarly generalize the Johnson result to get an algorithm for UG on HDXs via the steps outlined in Section~\ref{sec:abstraction}. To summarize, component 1 therein -- an SoS certificate of global hypercontractivity for HDXs -- can be obtained using~\cite{bafna2022hypercontractivity}, components 2 and 3 work for all graphs, and  component 4 is an iteration result that we can obtain, akin to the iteration result in~\cite{DBLP:conf/soda/BafnaHKL22}.
%The proof of this theorem requires an SoS certificate of Theorem~\ref{thm:expansion} and this turns out to be tedious to present. Such a certificate exists since the proof of~\Cref{thm:expansion} was exactly along the lines of KMMS' proof for Johnson graphs which we showed how to prove in the low-degree SoS proof system (Theorem~\ref{thm:structure-johnson}). So instead of reproving \Cref{thm:expansion} in the SoS proof system we reprove the weaker expansion result presented in Chapter~\ref{chap:hdx-tr} in the SoS system (Theorem~\ref{thm:hdx-expansion}) and use this result to give an algorithm for HD-walks albeit with weaker approximation guarantees that depend on the dimension of the complex:

%Recall that the parameters in the weaker structure theorem from Chapter~\ref{chap:hdx-tr} depended on $k$ too and therefore if one uses this instead of the stronger structure theorem from Chapter~\ref{chap:hdx-hypcont} the decay with $k$ in the algorithmic guarantee is unavoidable. But the proof of the above theorem highlights the major differences from the simpler Johnson case hence it is worth presenting. Johnson graphs are regular and have many symmetries and it might seem that our analysis depended on such properties. In this section by proving Theorem~\ref{thm:ug-hdx-weak} we show that our algorithmic framework is much more general and does not depend on such symmetry properties, since HD-walks could be over irregular sparse and highly asymmetric graphs.







